% القسم: المناهج
% الفصل: البراهين
% المبحث: قراءة البراهين

\documentclass[../../../include/open-logic-section]{subfiles}

\begin{document}

\olfileid{mth}{prf}{rea}
\olsection{قراءة البراهين}

قلّ أن تستوفي براهين الكتب والمقالات كل تفصيل أظهرناه في الأمثلة
المتقدمة. فكثيرًا ما يطوي المصنف ذكر تقسيم الحالات أو سلوك الطريق
غير المباشر، ولا ينص على التعريف الذي أعمله. فعلى قارئ البرهان أن
يستخرج تلك الوسائط بنفسه ليتم له الفهم. وفي هذا أيضًا رياضة نافعة على
إتقان خطوات البرهان المختلفة. ولننظر في مثال.

\begin{prop}[الامتصاص]
لكل مجموعتين $A$ و$B$،
\[
A \cap (A \cup B) = A.
\]
\end{prop}

\begin{proof}
إذا كان $z \in A \cap (A \cup B)$، فإن $z \in A$، ومن ثم
$A \cap (A \cup B) \subseteq A$. والآن افترض $z \in A$. عندئذ يكون
$z \in A \cup B$ أيضاً، ولذلك يكون $z \in A \cap (A \cup B)$ أيضاً.
\end{proof}

هذا البرهان لقانون الامتصاص بالغ الإيجاز: لم يسم التعريفات المستعملة،
ولم يقدم كل مطلوب بقول «علينا أن نثبت»، وما أشبه ذلك. فلنبسط مطويه.
القضية كلية في مجموعتين كيف اتفقتا، هما $A$ و$B$؛ فكلما ورد $A$ أو
$B$ كان متغيرًا لمجموعة اعتباطية. وما يثبته البرهان هو الشرط الكلي
لتطابق مجموعتين: كل !!{element} في الطرف الأيسر من المتطابقة
!!a{element} في الطرف الأيمن، والعكس.

\begin{quote}
«إذا كان $z \in A \cap (A \cup B)$، فإن $z \in A$، ومن ثم
$A \cap (A \cup B) \subseteq A$».
\end{quote}

هذا نصف البرهان الأول؛ فهو يأخذ $z$ اعتباطيًا !!a{element} من الطرف
الأيسر، ويثبت أنه !!a{element} أيضًا في الطرف الأيمن، أي يثبت
$A \cap (A \cup B) \subseteq A$. فإذا فرضنا
$z \in A \cap (A \cup B)$، وكان $z$ لا يكون !!{element} في تقاطع
مجموعتين إلا إذا كان !!{element} في كل واحدة منهما، نتج
$z \in A$ ومعه $z \in A \cup B$. وبوجه خاص يصدق $z \in A$، وهو
المطلوب. فإذا تم هذا النصف، تعين أن يكون باقي الكلام للنصف الثاني، أي
لإثبات
$A \subseteq A \cap (A \cup B)$.

\begin{quote}
«والآن افترض $z \in A$. عندئذ يكون $z \in A \cup B$ أيضاً، ولذلك يكون
$z \in A \cap (A \cup B)$ أيضاً».
\end{quote}

ومبدأ هذا النصف فرض $z \in A$، لأن المراد بيان أنه لكل~$z$: إذا كان
$z \in A$ كان $z \in A \cap (A \cup B)$. وإثبات
$z \in A \cap (A \cup B)$ يقتضي، بتعريف «$\cap$»، أمرين: (1)
$z \in A$، و(2) $z \in A \cup B$. أما (1) فهو نفس الفرض، فلا معنى
لإعادته. وأما (2)، فـ$z$ يكون !!a{element} في اتحاد مجموعتين إذا
وفقط إذا كان !!{element} في واحدة منهما على الأقل. ولدينا
$z \in A$، و$A \cup B$ هو اتحاد $A$ و$B$؛ فثبت
$z \in A \cup B$. وقد حصل (1) $z \in A$ و(2)
$z \in A \cup B$، ومن ثم يعطي تعريف «$\cap$» النتيجة
$z \in A \cap (A \cup B)$. وإنما طوى البرهان ذكر التعريفات لافتراضه
رسوخها عند القارئ. فمن لم ترسخ عنده فليرجع إليها، وليتحقق أيضًا من
أنه إذا كان $z \in A$ لزم $z \in A \cup B$، وأنه إذا اجتمع
$z \in A$ و$z \in A \cup B$ لزم $z \in A \cap (A \cup B)$.

وهذه صيغة أخرى للبرهان نفسه تُظهر ما طواه الأول:
\begin{proof}{}
[بحسب تعريف $=$ للمجموعات، لإثبات $A \cap (A \cup B) = A$ علينا أن
  نبين (أ) $A \cap (A \cup B) \subseteq A$ و(ب)
  $A \subseteq A \cap (A \cup B)$. (أ): بحسب تعريف $\subseteq$، علينا
  أن نبين أنه إذا كان $z \in A \cap (A \cup B)$ فإن $z \in A$.]
إذا كان $z \in A \cap (A \cup B)$، فإن $z \in A$ [لأنه بحسب تعريف
$\cap$، يكون $z \in A \cap (A \cup B)$ إذا وفقط إذا كان $z \in A$
و$z \in A \cup B$]، ومن ثم $A \cap (A \cup B) \subseteq A$.
[(ب): بحسب تعريف $\subseteq$، علينا أن نبين أنه إذا كان $z \in A$ فإن
$z \in A \cap (A \cup B)$.] افترض الآن [(1)] $z \in A$. عندئذ أيضاً
[(2)] $z \in A \cup B$ [لأن (1) يعطينا $z \in A$ أو $z \in B$،
وهو ما يعني، بحسب تعريف $\cup$، أن $z \in A \cup B$]؛ ومن ثم يكون
$z \in A \cap (A \cup B)$ أيضاً [لأن تعريف $\cap$ يقتضي
$z \in A$، أي (1)، و$z \in A \cup B$، أي (2)].
\end{proof}

\begin{prob}
ابسط البرهان الآتي لـ$A \cup (A \cap B) = A$، وبيّن أنماط الاستدلال
كلها، ووجه لزوم كل خطوة عن الفروض أو القضايا السابقة، ومواضع الحاجة
إلى كل تعريف.
\begin{proof}
  إذا كان $z \in A \cup (A \cap B)$ فإن $z \in A$ أو
  $z \in A \cap B$. وإذا كان $z \in A \cap B$ فإن $z \in A$. وكل
  $z \in A$ ينتمي أيضاً إلى $A \cup (A \cap B)$.
\end{proof}
\end{prob}

\end{document}
